%%%%%%%%%%%%%%%%%%%%%%%%%%%%%%%%%%%%%%%%%%%%%%%%%%%%%%%%%%%%%%%
% IMPORTANTE: compilar con XeLaTeX (fontspec lo requiere).      %
% Preámbulo de estilo para la presentación en Beamer (16:9),    %
% diseño académico sobrio para proyección en una sala mediana.  %
% Se importa desde main.tex con %%%%%%%%%%%%%%%%%%%%%%%%%%%%%%%%%%%%%%%%%%%%%%%%%%%%%%%%%%%%%%%
% IMPORTANTE: compilar con XeLaTeX (fontspec lo requiere).      %
% Preámbulo de estilo para la presentación en Beamer (16:9),    %
% diseño académico sobrio para proyección en una sala mediana.  %
% Se importa desde main.tex con %%%%%%%%%%%%%%%%%%%%%%%%%%%%%%%%%%%%%%%%%%%%%%%%%%%%%%%%%%%%%%%
% IMPORTANTE: compilar con XeLaTeX (fontspec lo requiere).      %
% Preámbulo de estilo para la presentación en Beamer (16:9),    %
% diseño académico sobrio para proyección en una sala mediana.  %
% Se importa desde main.tex con %%%%%%%%%%%%%%%%%%%%%%%%%%%%%%%%%%%%%%%%%%%%%%%%%%%%%%%%%%%%%%%
% IMPORTANTE: compilar con XeLaTeX (fontspec lo requiere).      %
% Preámbulo de estilo para la presentación en Beamer (16:9),    %
% diseño académico sobrio para proyección en una sala mediana.  %
% Se importa desde main.tex con \input{style}.                  %
%%%%%%%%%%%%%%%%%%%%%%%%%%%%%%%%%%%%%%%%%%%%%%%%%%%%%%%%%%%%%%%

%----------------------------- IDIOMA Y FUENTES ---------------------------
\usepackage[spanish,provide=*,es-noshorthands]{babel}
\selectlanguage{spanish}

\usepackage{fontspec}
% Fuente profesional opcional (requiere tenerla instalada). Por defecto se
% usa la fuente sans-serif estándar de XeLaTeX (Latin Modern Sans), que no
% requiere instalación adicional. Para activar otra, descomentar:
% \setsansfont{Lato}
% \setmainfont{Lato}

%----------------------------- PAQUETES BASE ---------------------------
\usepackage{graphicx}
\usepackage{booktabs}
\usepackage{tabularx}
\usepackage{array}
\usepackage{xcolor}
\usepackage{amsmath,amssymb}
\usepackage{siunitx}
\usepackage{tikz}
\usetikzlibrary{arrows.meta,positioning}
\usepackage{enumitem}

%----------------------------- PALETA (Okabe-Ito, igual que las figuras) ---
\definecolor{navy}{HTML}{1B2A4A}           % títulos
\definecolor{grafito}{HTML}{3C3C3C}        % texto
\definecolor{azulacento}{HTML}{0072B2}     % DE / Ridge / entrenamiento
\definecolor{naranjaacento}{HTML}{E69F00}  % M2 / validación
\definecolor{verdeacento}{HTML}{009E73}    % M3
\definecolor{grisreferencia}{HTML}{999999} % referencias
\definecolor{naranjaoscuro}{HTML}{D55E00}  % M3* cíclico / FR

%----------------------------- TEMA BASE (sin decoraciones) ---------------
\usetheme{default}
\usecolortheme{default}
\setbeamertemplate{navigation symbols}{}
\setbeamertemplate{blocks}[rounded][shadow=false]
\setbeamertemplate{itemize items}[circle]
\setbeamertemplate{enumerate items}[default]
\setbeamertemplate{headline}{} % sin barra superior

\setbeamercolor{background canvas}{bg=white}
\setbeamercolor{normal text}{fg=grafito,bg=white}
\setbeamercolor{structure}{fg=azulacento}
\setbeamercolor{title}{fg=navy}
\setbeamercolor{subtitle}{fg=grafito}
\setbeamercolor{frametitle}{fg=navy,bg=white}
\setbeamercolor{section in head/foot}{fg=grisreferencia,bg=white}
\setbeamercolor{page number in head/foot}{fg=grisreferencia,bg=white}
\setbeamercolor{itemize item}{fg=azulacento}
\setbeamercolor{itemize subitem}{fg=verdeacento}
\setbeamercolor{enumerate item}{fg=azulacento}
\setbeamercolor{alerted text}{fg=naranjaoscuro}
\setbeamercolor{block title}{fg=navy,bg=white}
\setbeamercolor{block body}{fg=grafito,bg=white}
\setbeamercolor{caption name}{fg=grisreferencia}

%----------------------------- TIPOGRAFÍA -----------------------------
% Tamaños relativos de Beamer (no \fontsize absolutos): la página de Beamer
% en aspectratio=169 es un lienzo pequeño (160mm x 90mm) que se amplifica al
% proyectarse, por lo que \fontsize{30}{36}pt literales (pensados para una
% hoja impresa normal) desbordan el frame. Se usan los comandos relativos
% estándar de LaTeX/Beamer, escalados desde el \\documentclass[11pt]{beamer}.
\setbeamerfont{title}{size=\Huge,series=\bfseries}
\setbeamerfont{subtitle}{size=\large}
\setbeamerfont{frametitle}{size=\LARGE,series=\bfseries}
\setbeamerfont{normal text}{size=\normalsize}
\setbeamerfont{itemize/enumerate body}{size=\normalsize}
\setbeamerfont{itemize/enumerate subbody}{size=\small}
\setbeamerfont{block title}{size=\small,series=\bfseries}
\setbeamerfont{block body}{size=\footnotesize}
\setbeamerfont{caption}{size=\footnotesize}
\setbeamerfont{footline}{size=\tiny}
\renewcommand{\arraystretch}{1.25}
% Listas compactas: el espaciado por defecto de itemize/enumerate en Beamer
% es generoso y, acumulado en varios ítems, desborda el cuerpo del frame.
\setlist[itemize]{itemsep=2pt,topsep=2pt,parsep=0pt}
\setlist[enumerate]{itemsep=2pt,topsep=2pt,parsep=0pt}

%----------------------------- MÁRGENES Y PIE DE PÁGINA -----------------
\setbeamersize{text margin left=1.1cm,text margin right=1.1cm}

\setbeamertemplate{frametitle}{
  \vspace{0.25cm}
  \usebeamerfont{frametitle}\usebeamercolor[fg]{frametitle}\insertframetitle\\[2pt]
  {\color{grisreferencia}\rule{\linewidth}{0.4pt}}
}

% Pie discreto: nombre de la sección a la izquierda, numeración a la derecha.
% Sin barras de color, sombras ni degradados.
\setbeamertemplate{footline}{
  \leavevmode
  \hbox{%
  \begin{beamercolorbox}[wd=.75\paperwidth,ht=2.4ex,dp=1.1ex,leftskip=1.1cm]{section in head/foot}
    \usebeamerfont{footline}\insertsectionhead
  \end{beamercolorbox}%
  \begin{beamercolorbox}[wd=.25\paperwidth,ht=2.4ex,dp=1.1ex,right,rightskip=1.1cm]{page number in head/foot}
    \usebeamerfont{footline}\insertframenumber{} / \inserttotalframenumber
  \end{beamercolorbox}}%
}

%----------------------------- NOTA (RECURSO PENDIENTE) -----------------
% Para marcar, sin inventar datos, un recurso visual que aún no se adaptó.
\newcommand{\recursopendiente}[1]{%
  {\small\color{grisreferencia}\itshape #1}%
}

%----------------------------- AJUSTE DE TABLAS (SOLO ACHICAR) -----------
% \input{} de una tabla y la reduce, de forma uniforme (mismo factor en
% ancho y alto), solo lo necesario para que quepa en el ancho disponible
% (\linewidth) y en una altura máxima (por defecto, 0.42\textheight, holgada
% para no chocar con el resto del contenido del frame). Las tablas angostas
% o bajas no se agrandan (evita estirar el texto y desbordar el frame).
\newsavebox{\tablafitbox}
\newcommand{\tablafit}[2][0.42\textheight]{%
  \sbox{\tablafitbox}{\input{#2}}%
  \pgfmathsetmacro{\tablafitsw}{\linewidth/\wd\tablafitbox}%
  \pgfmathsetmacro{\tablafitsh}{(#1)/\dimexpr\ht\tablafitbox+\dp\tablafitbox\relax}%
  \pgfmathsetmacro{\tablafits}{min(1,\tablafitsw,\tablafitsh)}%
  \centerline{\resizebox{\tablafits\wd\tablafitbox}{!}{\usebox{\tablafitbox}}}%
}
.                  %
%%%%%%%%%%%%%%%%%%%%%%%%%%%%%%%%%%%%%%%%%%%%%%%%%%%%%%%%%%%%%%%

%----------------------------- IDIOMA Y FUENTES ---------------------------
\usepackage[spanish,provide=*,es-noshorthands]{babel}
\selectlanguage{spanish}

\usepackage{fontspec}
% Fuente profesional opcional (requiere tenerla instalada). Por defecto se
% usa la fuente sans-serif estándar de XeLaTeX (Latin Modern Sans), que no
% requiere instalación adicional. Para activar otra, descomentar:
% \setsansfont{Lato}
% \setmainfont{Lato}

%----------------------------- PAQUETES BASE ---------------------------
\usepackage{graphicx}
\usepackage{booktabs}
\usepackage{tabularx}
\usepackage{array}
\usepackage{xcolor}
\usepackage{amsmath,amssymb}
\usepackage{siunitx}
\usepackage{tikz}
\usetikzlibrary{arrows.meta,positioning}
\usepackage{enumitem}

%----------------------------- PALETA (Okabe-Ito, igual que las figuras) ---
\definecolor{navy}{HTML}{1B2A4A}           % títulos
\definecolor{grafito}{HTML}{3C3C3C}        % texto
\definecolor{azulacento}{HTML}{0072B2}     % DE / Ridge / entrenamiento
\definecolor{naranjaacento}{HTML}{E69F00}  % M2 / validación
\definecolor{verdeacento}{HTML}{009E73}    % M3
\definecolor{grisreferencia}{HTML}{999999} % referencias
\definecolor{naranjaoscuro}{HTML}{D55E00}  % M3* cíclico / FR

%----------------------------- TEMA BASE (sin decoraciones) ---------------
\usetheme{default}
\usecolortheme{default}
\setbeamertemplate{navigation symbols}{}
\setbeamertemplate{blocks}[rounded][shadow=false]
\setbeamertemplate{itemize items}[circle]
\setbeamertemplate{enumerate items}[default]
\setbeamertemplate{headline}{} % sin barra superior

\setbeamercolor{background canvas}{bg=white}
\setbeamercolor{normal text}{fg=grafito,bg=white}
\setbeamercolor{structure}{fg=azulacento}
\setbeamercolor{title}{fg=navy}
\setbeamercolor{subtitle}{fg=grafito}
\setbeamercolor{frametitle}{fg=navy,bg=white}
\setbeamercolor{section in head/foot}{fg=grisreferencia,bg=white}
\setbeamercolor{page number in head/foot}{fg=grisreferencia,bg=white}
\setbeamercolor{itemize item}{fg=azulacento}
\setbeamercolor{itemize subitem}{fg=verdeacento}
\setbeamercolor{enumerate item}{fg=azulacento}
\setbeamercolor{alerted text}{fg=naranjaoscuro}
\setbeamercolor{block title}{fg=navy,bg=white}
\setbeamercolor{block body}{fg=grafito,bg=white}
\setbeamercolor{caption name}{fg=grisreferencia}

%----------------------------- TIPOGRAFÍA -----------------------------
% Tamaños relativos de Beamer (no \fontsize absolutos): la página de Beamer
% en aspectratio=169 es un lienzo pequeño (160mm x 90mm) que se amplifica al
% proyectarse, por lo que \fontsize{30}{36}pt literales (pensados para una
% hoja impresa normal) desbordan el frame. Se usan los comandos relativos
% estándar de LaTeX/Beamer, escalados desde el \\documentclass[11pt]{beamer}.
\setbeamerfont{title}{size=\Huge,series=\bfseries}
\setbeamerfont{subtitle}{size=\large}
\setbeamerfont{frametitle}{size=\LARGE,series=\bfseries}
\setbeamerfont{normal text}{size=\normalsize}
\setbeamerfont{itemize/enumerate body}{size=\normalsize}
\setbeamerfont{itemize/enumerate subbody}{size=\small}
\setbeamerfont{block title}{size=\small,series=\bfseries}
\setbeamerfont{block body}{size=\footnotesize}
\setbeamerfont{caption}{size=\footnotesize}
\setbeamerfont{footline}{size=\tiny}
\renewcommand{\arraystretch}{1.25}
% Listas compactas: el espaciado por defecto de itemize/enumerate en Beamer
% es generoso y, acumulado en varios ítems, desborda el cuerpo del frame.
\setlist[itemize]{itemsep=2pt,topsep=2pt,parsep=0pt}
\setlist[enumerate]{itemsep=2pt,topsep=2pt,parsep=0pt}

%----------------------------- MÁRGENES Y PIE DE PÁGINA -----------------
\setbeamersize{text margin left=1.1cm,text margin right=1.1cm}

\setbeamertemplate{frametitle}{
  \vspace{0.25cm}
  \usebeamerfont{frametitle}\usebeamercolor[fg]{frametitle}\insertframetitle\\[2pt]
  {\color{grisreferencia}\rule{\linewidth}{0.4pt}}
}

% Pie discreto: nombre de la sección a la izquierda, numeración a la derecha.
% Sin barras de color, sombras ni degradados.
\setbeamertemplate{footline}{
  \leavevmode
  \hbox{%
  \begin{beamercolorbox}[wd=.75\paperwidth,ht=2.4ex,dp=1.1ex,leftskip=1.1cm]{section in head/foot}
    \usebeamerfont{footline}\insertsectionhead
  \end{beamercolorbox}%
  \begin{beamercolorbox}[wd=.25\paperwidth,ht=2.4ex,dp=1.1ex,right,rightskip=1.1cm]{page number in head/foot}
    \usebeamerfont{footline}\insertframenumber{} / \inserttotalframenumber
  \end{beamercolorbox}}%
}

%----------------------------- NOTA (RECURSO PENDIENTE) -----------------
% Para marcar, sin inventar datos, un recurso visual que aún no se adaptó.
\newcommand{\recursopendiente}[1]{%
  {\small\color{grisreferencia}\itshape #1}%
}

%----------------------------- AJUSTE DE TABLAS (SOLO ACHICAR) -----------
% \input{} de una tabla y la reduce, de forma uniforme (mismo factor en
% ancho y alto), solo lo necesario para que quepa en el ancho disponible
% (\linewidth) y en una altura máxima (por defecto, 0.42\textheight, holgada
% para no chocar con el resto del contenido del frame). Las tablas angostas
% o bajas no se agrandan (evita estirar el texto y desbordar el frame).
\newsavebox{\tablafitbox}
\newcommand{\tablafit}[2][0.42\textheight]{%
  \sbox{\tablafitbox}{\input{#2}}%
  \pgfmathsetmacro{\tablafitsw}{\linewidth/\wd\tablafitbox}%
  \pgfmathsetmacro{\tablafitsh}{(#1)/\dimexpr\ht\tablafitbox+\dp\tablafitbox\relax}%
  \pgfmathsetmacro{\tablafits}{min(1,\tablafitsw,\tablafitsh)}%
  \centerline{\resizebox{\tablafits\wd\tablafitbox}{!}{\usebox{\tablafitbox}}}%
}
.                  %
%%%%%%%%%%%%%%%%%%%%%%%%%%%%%%%%%%%%%%%%%%%%%%%%%%%%%%%%%%%%%%%

%----------------------------- IDIOMA Y FUENTES ---------------------------
\usepackage[spanish,provide=*,es-noshorthands]{babel}
\selectlanguage{spanish}

\usepackage{fontspec}
% Fuente profesional opcional (requiere tenerla instalada). Por defecto se
% usa la fuente sans-serif estándar de XeLaTeX (Latin Modern Sans), que no
% requiere instalación adicional. Para activar otra, descomentar:
% \setsansfont{Lato}
% \setmainfont{Lato}

%----------------------------- PAQUETES BASE ---------------------------
\usepackage{graphicx}
\usepackage{booktabs}
\usepackage{tabularx}
\usepackage{array}
\usepackage{xcolor}
\usepackage{amsmath,amssymb}
\usepackage{siunitx}
\usepackage{tikz}
\usetikzlibrary{arrows.meta,positioning}
\usepackage{enumitem}

%----------------------------- PALETA (Okabe-Ito, igual que las figuras) ---
\definecolor{navy}{HTML}{1B2A4A}           % títulos
\definecolor{grafito}{HTML}{3C3C3C}        % texto
\definecolor{azulacento}{HTML}{0072B2}     % DE / Ridge / entrenamiento
\definecolor{naranjaacento}{HTML}{E69F00}  % M2 / validación
\definecolor{verdeacento}{HTML}{009E73}    % M3
\definecolor{grisreferencia}{HTML}{999999} % referencias
\definecolor{naranjaoscuro}{HTML}{D55E00}  % M3* cíclico / FR

%----------------------------- TEMA BASE (sin decoraciones) ---------------
\usetheme{default}
\usecolortheme{default}
\setbeamertemplate{navigation symbols}{}
\setbeamertemplate{blocks}[rounded][shadow=false]
\setbeamertemplate{itemize items}[circle]
\setbeamertemplate{enumerate items}[default]
\setbeamertemplate{headline}{} % sin barra superior

\setbeamercolor{background canvas}{bg=white}
\setbeamercolor{normal text}{fg=grafito,bg=white}
\setbeamercolor{structure}{fg=azulacento}
\setbeamercolor{title}{fg=navy}
\setbeamercolor{subtitle}{fg=grafito}
\setbeamercolor{frametitle}{fg=navy,bg=white}
\setbeamercolor{section in head/foot}{fg=grisreferencia,bg=white}
\setbeamercolor{page number in head/foot}{fg=grisreferencia,bg=white}
\setbeamercolor{itemize item}{fg=azulacento}
\setbeamercolor{itemize subitem}{fg=verdeacento}
\setbeamercolor{enumerate item}{fg=azulacento}
\setbeamercolor{alerted text}{fg=naranjaoscuro}
\setbeamercolor{block title}{fg=navy,bg=white}
\setbeamercolor{block body}{fg=grafito,bg=white}
\setbeamercolor{caption name}{fg=grisreferencia}

%----------------------------- TIPOGRAFÍA -----------------------------
% Tamaños relativos de Beamer (no \fontsize absolutos): la página de Beamer
% en aspectratio=169 es un lienzo pequeño (160mm x 90mm) que se amplifica al
% proyectarse, por lo que \fontsize{30}{36}pt literales (pensados para una
% hoja impresa normal) desbordan el frame. Se usan los comandos relativos
% estándar de LaTeX/Beamer, escalados desde el \\documentclass[11pt]{beamer}.
\setbeamerfont{title}{size=\Huge,series=\bfseries}
\setbeamerfont{subtitle}{size=\large}
\setbeamerfont{frametitle}{size=\LARGE,series=\bfseries}
\setbeamerfont{normal text}{size=\normalsize}
\setbeamerfont{itemize/enumerate body}{size=\normalsize}
\setbeamerfont{itemize/enumerate subbody}{size=\small}
\setbeamerfont{block title}{size=\small,series=\bfseries}
\setbeamerfont{block body}{size=\footnotesize}
\setbeamerfont{caption}{size=\footnotesize}
\setbeamerfont{footline}{size=\tiny}
\renewcommand{\arraystretch}{1.25}
% Listas compactas: el espaciado por defecto de itemize/enumerate en Beamer
% es generoso y, acumulado en varios ítems, desborda el cuerpo del frame.
\setlist[itemize]{itemsep=2pt,topsep=2pt,parsep=0pt}
\setlist[enumerate]{itemsep=2pt,topsep=2pt,parsep=0pt}

%----------------------------- MÁRGENES Y PIE DE PÁGINA -----------------
\setbeamersize{text margin left=1.1cm,text margin right=1.1cm}

\setbeamertemplate{frametitle}{
  \vspace{0.25cm}
  \usebeamerfont{frametitle}\usebeamercolor[fg]{frametitle}\insertframetitle\\[2pt]
  {\color{grisreferencia}\rule{\linewidth}{0.4pt}}
}

% Pie discreto: nombre de la sección a la izquierda, numeración a la derecha.
% Sin barras de color, sombras ni degradados.
\setbeamertemplate{footline}{
  \leavevmode
  \hbox{%
  \begin{beamercolorbox}[wd=.75\paperwidth,ht=2.4ex,dp=1.1ex,leftskip=1.1cm]{section in head/foot}
    \usebeamerfont{footline}\insertsectionhead
  \end{beamercolorbox}%
  \begin{beamercolorbox}[wd=.25\paperwidth,ht=2.4ex,dp=1.1ex,right,rightskip=1.1cm]{page number in head/foot}
    \usebeamerfont{footline}\insertframenumber{} / \inserttotalframenumber
  \end{beamercolorbox}}%
}

%----------------------------- NOTA (RECURSO PENDIENTE) -----------------
% Para marcar, sin inventar datos, un recurso visual que aún no se adaptó.
\newcommand{\recursopendiente}[1]{%
  {\small\color{grisreferencia}\itshape #1}%
}

%----------------------------- AJUSTE DE TABLAS (SOLO ACHICAR) -----------
% \input{} de una tabla y la reduce, de forma uniforme (mismo factor en
% ancho y alto), solo lo necesario para que quepa en el ancho disponible
% (\linewidth) y en una altura máxima (por defecto, 0.42\textheight, holgada
% para no chocar con el resto del contenido del frame). Las tablas angostas
% o bajas no se agrandan (evita estirar el texto y desbordar el frame).
\newsavebox{\tablafitbox}
\newcommand{\tablafit}[2][0.42\textheight]{%
  \sbox{\tablafitbox}{\input{#2}}%
  \pgfmathsetmacro{\tablafitsw}{\linewidth/\wd\tablafitbox}%
  \pgfmathsetmacro{\tablafitsh}{(#1)/\dimexpr\ht\tablafitbox+\dp\tablafitbox\relax}%
  \pgfmathsetmacro{\tablafits}{min(1,\tablafitsw,\tablafitsh)}%
  \centerline{\resizebox{\tablafits\wd\tablafitbox}{!}{\usebox{\tablafitbox}}}%
}
.                  %
%%%%%%%%%%%%%%%%%%%%%%%%%%%%%%%%%%%%%%%%%%%%%%%%%%%%%%%%%%%%%%%

%----------------------------- IDIOMA Y FUENTES ---------------------------
\usepackage[spanish,provide=*,es-noshorthands]{babel}
\selectlanguage{spanish}

\usepackage{fontspec}
% Fuente profesional opcional (requiere tenerla instalada). Por defecto se
% usa la fuente sans-serif estándar de XeLaTeX (Latin Modern Sans), que no
% requiere instalación adicional. Para activar otra, descomentar:
% \setsansfont{Lato}
% \setmainfont{Lato}

%----------------------------- PAQUETES BASE ---------------------------
\usepackage{graphicx}
\usepackage{booktabs}
\usepackage{tabularx}
\usepackage{array}
\usepackage{xcolor}
\usepackage{amsmath,amssymb}
\usepackage{siunitx}
\usepackage{tikz}
\usetikzlibrary{arrows.meta,positioning}
\usepackage{enumitem}

%----------------------------- PALETA (Okabe-Ito, igual que las figuras) ---
\definecolor{navy}{HTML}{1B2A4A}           % títulos
\definecolor{grafito}{HTML}{3C3C3C}        % texto
\definecolor{azulacento}{HTML}{0072B2}     % DE / Ridge / entrenamiento
\definecolor{naranjaacento}{HTML}{E69F00}  % M2 / validación
\definecolor{verdeacento}{HTML}{009E73}    % M3
\definecolor{grisreferencia}{HTML}{999999} % referencias
\definecolor{naranjaoscuro}{HTML}{D55E00}  % M3* cíclico / FR

%----------------------------- TEMA BASE (sin decoraciones) ---------------
\usetheme{default}
\usecolortheme{default}
\setbeamertemplate{navigation symbols}{}
\setbeamertemplate{blocks}[rounded][shadow=false]
\setbeamertemplate{itemize items}[circle]
\setbeamertemplate{enumerate items}[default]
\setbeamertemplate{headline}{} % sin barra superior

\setbeamercolor{background canvas}{bg=white}
\setbeamercolor{normal text}{fg=grafito,bg=white}
\setbeamercolor{structure}{fg=azulacento}
\setbeamercolor{title}{fg=navy}
\setbeamercolor{subtitle}{fg=grafito}
\setbeamercolor{frametitle}{fg=navy,bg=white}
\setbeamercolor{section in head/foot}{fg=grisreferencia,bg=white}
\setbeamercolor{page number in head/foot}{fg=grisreferencia,bg=white}
\setbeamercolor{itemize item}{fg=azulacento}
\setbeamercolor{itemize subitem}{fg=verdeacento}
\setbeamercolor{enumerate item}{fg=azulacento}
\setbeamercolor{alerted text}{fg=naranjaoscuro}
\setbeamercolor{block title}{fg=navy,bg=white}
\setbeamercolor{block body}{fg=grafito,bg=white}
\setbeamercolor{caption name}{fg=grisreferencia}

%----------------------------- TIPOGRAFÍA -----------------------------
% Tamaños relativos de Beamer (no \fontsize absolutos): la página de Beamer
% en aspectratio=169 es un lienzo pequeño (160mm x 90mm) que se amplifica al
% proyectarse, por lo que \fontsize{30}{36}pt literales (pensados para una
% hoja impresa normal) desbordan el frame. Se usan los comandos relativos
% estándar de LaTeX/Beamer, escalados desde el \\documentclass[11pt]{beamer}.
\setbeamerfont{title}{size=\Huge,series=\bfseries}
\setbeamerfont{subtitle}{size=\large}
\setbeamerfont{frametitle}{size=\LARGE,series=\bfseries}
\setbeamerfont{normal text}{size=\normalsize}
\setbeamerfont{itemize/enumerate body}{size=\normalsize}
\setbeamerfont{itemize/enumerate subbody}{size=\small}
\setbeamerfont{block title}{size=\small,series=\bfseries}
\setbeamerfont{block body}{size=\footnotesize}
\setbeamerfont{caption}{size=\footnotesize}
\setbeamerfont{footline}{size=\tiny}
\renewcommand{\arraystretch}{1.25}
% Listas compactas: el espaciado por defecto de itemize/enumerate en Beamer
% es generoso y, acumulado en varios ítems, desborda el cuerpo del frame.
\setlist[itemize]{itemsep=2pt,topsep=2pt,parsep=0pt}
\setlist[enumerate]{itemsep=2pt,topsep=2pt,parsep=0pt}

%----------------------------- MÁRGENES Y PIE DE PÁGINA -----------------
\setbeamersize{text margin left=1.1cm,text margin right=1.1cm}

\setbeamertemplate{frametitle}{
  \vspace{0.25cm}
  \usebeamerfont{frametitle}\usebeamercolor[fg]{frametitle}\insertframetitle\\[2pt]
  {\color{grisreferencia}\rule{\linewidth}{0.4pt}}
}

% Pie discreto: nombre de la sección a la izquierda, numeración a la derecha.
% Sin barras de color, sombras ni degradados.
\setbeamertemplate{footline}{
  \leavevmode
  \hbox{%
  \begin{beamercolorbox}[wd=.75\paperwidth,ht=2.4ex,dp=1.1ex,leftskip=1.1cm]{section in head/foot}
    \usebeamerfont{footline}\insertsectionhead
  \end{beamercolorbox}%
  \begin{beamercolorbox}[wd=.25\paperwidth,ht=2.4ex,dp=1.1ex,right,rightskip=1.1cm]{page number in head/foot}
    \usebeamerfont{footline}\insertframenumber{} / \inserttotalframenumber
  \end{beamercolorbox}}%
}

%----------------------------- NOTA (RECURSO PENDIENTE) -----------------
% Para marcar, sin inventar datos, un recurso visual que aún no se adaptó.
\newcommand{\recursopendiente}[1]{%
  {\small\color{grisreferencia}\itshape #1}%
}

%----------------------------- AJUSTE DE TABLAS (SOLO ACHICAR) -----------
% \input{} de una tabla y la reduce, de forma uniforme (mismo factor en
% ancho y alto), solo lo necesario para que quepa en el ancho disponible
% (\linewidth) y en una altura máxima (por defecto, 0.42\textheight, holgada
% para no chocar con el resto del contenido del frame). Las tablas angostas
% o bajas no se agrandan (evita estirar el texto y desbordar el frame).
\newsavebox{\tablafitbox}
\newcommand{\tablafit}[2][0.42\textheight]{%
  \sbox{\tablafitbox}{\input{#2}}%
  \pgfmathsetmacro{\tablafitsw}{\linewidth/\wd\tablafitbox}%
  \pgfmathsetmacro{\tablafitsh}{(#1)/\dimexpr\ht\tablafitbox+\dp\tablafitbox\relax}%
  \pgfmathsetmacro{\tablafits}{min(1,\tablafitsw,\tablafitsh)}%
  \centerline{\resizebox{\tablafits\wd\tablafitbox}{!}{\usebox{\tablafitbox}}}%
}
